\section{Framework flexible corollaries}
\label{sec:framework-flexible-corollaries}

This note preserves framework variants that are useful for bookkeeping and
future construction lines but are not needed in the compiled main spine.  The
main paper keeps the first-moment theorem and the basic envelope corollary.

\subsection{Input--output enumerator viewpoint}

When the inner code \(\Ctwo\) is deterministic and linear, one can define the
slice enumerator
\[
\enum^{\Ctwo}_{w,h}
\;:=\;
\bigl|\{u\in\{0,1\}^n:\ \wt(u)=w,\ \wt(\Ctwo(u))=h\}\bigr|.
\]
Then
\[
\Pr[\wt(\Ctwo(U))=h\mid \wt(U)=w]
=
\frac{\enum^{\Ctwo}_{w,h}}{\binom{n}{w}},
\qquad
p_w(\delta)
=
\sum_{h\le \delta n}
\frac{\enum^{\Ctwo}_{w,h}}{\binom{n}{w}} .
\]
The accumulator warmup uses this viewpoint directly; the dense and full-split
inners are analyzed through direct slice-to-tail bounds.

\subsection{Piecewise evaluation}

\begin{corollary}[Piecewise evaluation of the framework sum]
\label{cor:framework-piecewise}
Under the hypotheses of Corollary~\ref{cor:framework-envelope}, let
\[
\{w_0,\dots,n\} \;=\; I_1 \sqcup I_2 \sqcup \cdots \sqcup I_J
\]
be any partition into weight ranges.
Suppose that for each \(j\in[J]\), and each \(w\in I_j\), we have
\[
\EE[\enum_{\Cone,w}] \;\le\; A^{(j)}_w,
\qquad
p_w(\delta)\;\le\; B^{(j)}_w.
\]
Then
\[
\Pr[d_{\min}(\Csc)\le \delta n]
\;\le\;
\varepsilon_{\mathrm{out}}(n)
\;+\;
\sum_{j=1}^{J}\ \sum_{w\in I_j} A^{(j)}_w\,B^{(j)}_w.
\]
\end{corollary}

This is only bookkeeping: low, intermediate, and linear weight regimes enter
through whichever outer and inner envelopes are most effective on those ranges.

\subsection{Single-envelope specialization}

\begin{corollary}[Single-envelope specialization]
\label{cor:linear-dist-from-interfaces}
Assume the outer interface (Assumption~\ref{ass:outer-interface}) and the inner
interface (Assumption~\ref{ass:inner-interface}) hold with parameters
\((\alpha,c,\beta)\) and \((a,\rho)\), respectively.  If \(\beta\rho<1\) and
\begin{equation}
\label{eq:poly-beating-condition}
\alpha\cdot \log\Big(\frac{1}{\beta\rho}\Big) \;>\; a+c,
\end{equation}
then
\[
\Pr\bigl[d_{\min}(\Csc)\le \delta n\bigr] \;=\; o(1)
\qquad\text{as }n\to\infty.
\]
In particular, \(d_{\min}(\Csc)=\Omega(n)\) with high probability.
\end{corollary}

\begin{proof}
Apply Corollary~\ref{cor:framework-envelope} with
\[
w_0:=\lceil \alpha\log n\rceil,
\qquad
\varepsilon_{\mathrm{out}}(n):=\Pr[d_{\min}(\Cone)<\alpha\log n]=o(1),
\]
and envelopes
\[
A_w:=n^c\beta^w,
\qquad
B_w:=n^a\rho^w.
\]
This gives
\[
\Pr[d_{\min}(\Csc)\le \delta n]
\;\le\;
o(1)\;+\;n^{a+c}\sum_{w=w_0}^{n} (\beta\rho)^w.
\]
Since \(\beta\rho<1\),
\[
\sum_{w=w_0}^{n}(\beta\rho)^w
\;\le\;
\frac{(\beta\rho)^{w_0}}{1-\beta\rho}
\;\le\;
O(1)\cdot n^{-\alpha\log(1/(\beta\rho))}.
\]
The exponent is negative exactly under
\eqref{eq:poly-beating-condition}.
\end{proof}
